\uuid{f6Cb}
\titre{Équation droite/plan}
\auteur{bodin ; hanani ; mzari}
\organisation{exo7}
\licence{CC BY-NC-SA 4.0}
\datecreate{2020-01-20}
\niveau{L1}
\module{Géométrie}
\chapitre{Géométrie de l'espace}
\sousChapitre{Droites de l'espace}
\difficulte{1}
\theme{Droites de l'espace, équation droite, plan}
\exo7qcmid{lille-1.geometrie-espace.007}

\contenu{
\texte{
Pour ces questions, l'espace est muni d'un repère orthonormé direct  $(O, \vec{i}, \vec{j}, \vec{k})$.
}

\begin{question}
\qclassification{141.04}
\qtags{motcle=équation droite/plan}
Soit $D$ la droite passant par le point $A(2,-1,1)$ et dirigée par le vecteur $\vec{u}(-1,1,0)$. Quelles sont les assertions vraies ?
\begin{answers}

    \good {Une représentation paramétrique de $D$ est :
 $$\left\{\begin{array}{ccl}x&=&2-t\\y&=&-1+t\\ z&=&1,\quad (t\in \Rr) \end{array}\right.$$}

    \bad {Une représentation paramétrique de $D$ est :
 $$\left\{\begin{array}{ccl}x&=&2-t\\y&=&-1+t\\ z&=&-t,\quad (t\in \Rr)\end{array}\right.$$}

   \good{Une représentation cartésienne de $D$ est :
  $$\left\{\begin{array}{ccl}x+y&=&1\\z&=&1 \end{array}\right.$$}

    \bad{Une représentation cartésienne de $D$ est :
  $$\left\{\begin{array}{ccl}x+y&=&0\\z&=&1 \end{array}\right.$$}

\end{answers}
\begin{explanations}
On peut trouver une représentation cartésienne, à partir d'une représentation   paramétrique en éliminant le paramètre.
\end{explanations}
\end{question}
}
